% GENERATED FILE. Edit canonical lessons, then run npm run generate:books.
% canonical-source-hash: fnv1a64:5103e9f2c3eba339
% canonical-lessons: JA-C457-shuushoku-suru, JA-C457-mensetsu, JA-C457-rirekisho, JA-C457-oubo-suru, JA-C457-boshuu

\chapter{Looking for a Job}
\label{ch:ja-looking-for-a-job}
\hlchaptermodality{457}

\begin{chapteropening}
\textbf{By the end of this chapter:} \emph{I can talk about applying for a job.}

Complete the last lesson of chapter 457: I can talk about applying for a job.
\end{chapteropening}

\section[\texorpdfstring{\ja{しゅうしょくする} (shūshoku suru)}{しゅうしょくする (shūshoku suru)}]{\ja{しゅうしょくする} (shūshoku suru) — to find a job}

\label{lesson:JA-C457-shuushoku-suru}



\begin{quote}
Before the new one: say the Japanese for to hire, then the Japanese for a workplace.
\end{quote}

\subsection*{You'll want to know: \ja{しゅうしょくする}}
\textbf{\ja{しゅうしょくする}} — \emph{shūshoku suru} — \textquotedblleft{}to find a job\textquotedblright{}.

\begin{quote}
\textbf{\ja{あねは} \ja{ぎんこうに} \ja{しゅうしょくしました。}} — \emph{ane wa ginkō ni shūshoku shimashita} — \textquotedblleft{}My elder sister got a job at a bank.\textquotedblright{}
\end{quote}

Polite \textbf{\ja{しゅうしょくします}}, \emph{shūshokushimasu}; a \ja{する} verb. Past (polite): \textbf{\ja{あにが} \ja{かいしゃに} \ja{しゅうしょくしました。}}, \emph{ani ga kaisha ni shūshoku shimashita}.

\begin{cousinweb}
\textbf{\ja{しゅうしょくする}} (\emph{shūshoku suru}) is Sino-Japanese (kango) \emph{shūshoku} with the native verb \emph{suru}, \textquotedblleft{}to do\textquotedblright{}: two Chinese characters read with their on'yomi, readings borrowed from Chinese. They are \emph{shū}, \textquotedblleft{}settle\textquotedblright{} (Middle Chinese \emph{dzjuwH}) and \emph{shoku}, \textquotedblleft{}employment\textquotedblright{} (\emph{tsyik}) (KANJIDIC2; Middle Chinese in Baxter's transcription).
\end{cousinweb}

\begin{grammarlens}[title={What you've built}]
One more word for describing what you lived and felt.
\end{grammarlens}

\subsection*{Your turn}
\begin{itemize}
\raggedright
  \item \emph{Say it:} \emph{shūshoku suru}
  \item \emph{Say it:} \emph{shūshoku suru}, once more
  \item \emph{Say it:} \emph{ane wa ginkō ni shūshoku shimashita} — My elder sister got a job at a bank.
  \item \emph{Recall:} say the Japanese for to hire, then the Japanese for a workplace, then say \emph{shūshoku suru} again
\end{itemize}

\begin{tcolorbox}[breakable,colback=teal!4,colframe=teal!35!black,title={Before you move on}]
What does \ja{しゅうしょくする} mean? (\textquotedblleft{}To find a job\textquotedblright{}.) Say it once more.
\end{tcolorbox}
\section[\texorpdfstring{\ja{めんせつ} (mensetsu)}{めんせつ (mensetsu)}]{\ja{めんせつ} (mensetsu) — a job interview; an interview}

\label{lesson:JA-C457-mensetsu}



\begin{quote}
Before the new one: say the Japanese for a workplace, then the Japanese for to find a job.
\end{quote}

\subsection*{You'll want to know: \ja{めんせつ}}
\textbf{\ja{めんせつ}} — \emph{mensetsu} — \textquotedblleft{}a job interview; an interview\textquotedblright{}.

\begin{quote}
\textbf{\ja{あした} \ja{かいしゃで} \ja{めんせつが} \ja{あります。}} — \emph{ashita kaisha de mensetsu ga arimasu} — \textquotedblleft{}I have an interview at a company tomorrow.\textquotedblright{}
\end{quote}

\begin{cousinweb}
\textbf{\ja{めんせつ}} (\emph{mensetsu}) is Sino-Japanese (kango): two Chinese characters read with their on'yomi, readings borrowed from Chinese. They are \emph{men}, \textquotedblleft{}face\textquotedblright{} (Middle Chinese \emph{mjienH}) and \emph{setsu}, \textquotedblleft{}contact\textquotedblright{} (\emph{tsjep}) (KANJIDIC2; Middle Chinese in Baxter's transcription).
\end{cousinweb}

\begin{grammarlens}[title={What you've built}]
One more word for describing what you lived and felt.
\end{grammarlens}

\subsection*{Your turn}
\begin{itemize}
\raggedright
  \item \emph{Say it:} \emph{mensetsu}
  \item \emph{Say it:} \emph{mensetsu}, once more
  \item \emph{Say it:} \emph{ashita kaisha de mensetsu ga arimasu} — I have an interview at a company tomorrow.
  \item \emph{Recall:} say the Japanese for a workplace, then the Japanese for to find a job, then say \emph{mensetsu} again
\end{itemize}

\begin{tcolorbox}[breakable,colback=teal!4,colframe=teal!35!black,title={Before you move on}]
What does \ja{めんせつ} mean? (\textquotedblleft{}A job interview; an interview\textquotedblright{}.) Say it once more.
\end{tcolorbox}
\section[\texorpdfstring{\ja{りれきしょ} (rirekisho)}{りれきしょ (rirekisho)}]{\ja{りれきしょ} (rirekisho) — a curriculum vitae}

\label{lesson:JA-C457-rirekisho}



\begin{quote}
Before the new one: say the Japanese for to find a job, then the Japanese for a job interview.
\end{quote}

\subsection*{You'll want to know: \ja{りれきしょ}}
\textbf{\ja{りれきしょ}} — \emph{rirekisho} — \textquotedblleft{}a curriculum vitae\textquotedblright{}.

\begin{quote}
\textbf{\ja{りれきしょを} \ja{かいしゃに} \ja{おくりました。}} — \emph{rirekisho o kaisha ni okurimashita} — \textquotedblleft{}I sent my curriculum vitae to the company.\textquotedblright{}
\end{quote}

\begin{cousinweb}
\textbf{\ja{りれきしょ}} (\emph{rirekisho}) is Sino-Japanese (kango): three Chinese characters read with their on'yomi, readings borrowed from Chinese. They are \emph{ri}, \textquotedblleft{}perform\textquotedblright{} (Middle Chinese \emph{lijX}), \emph{reki}, \textquotedblleft{}curriculum\textquotedblright{} (\emph{lek}), and \emph{sho}, \textquotedblleft{}write\textquotedblright{} (\emph{syo}) (KANJIDIC2; Middle Chinese in Baxter's transcription).
\end{cousinweb}

\begin{grammarlens}[title={What you've built}]
One more word for describing what you lived and felt.
\end{grammarlens}

\subsection*{Your turn}
\begin{itemize}
\raggedright
  \item \emph{Say it:} \emph{rirekisho}
  \item \emph{Say it:} \emph{rirekisho}, once more
  \item \emph{Say it:} \emph{rirekisho o kaisha ni okurimashita} — I sent my curriculum vitae to the company.
  \item \emph{Recall:} say the Japanese for to find a job, then the Japanese for a job interview, then say \emph{rirekisho} again
\end{itemize}

\begin{tcolorbox}[breakable,colback=teal!4,colframe=teal!35!black,title={Before you move on}]
What does \ja{りれきしょ} mean? (\textquotedblleft{}A curriculum vitae\textquotedblright{}.) Say it once more.
\end{tcolorbox}
\section[\texorpdfstring{\ja{おうぼする} (ōbo suru)}{おうぼする (ōbo suru)}]{\ja{おうぼする} (ōbo suru) — to apply (for a job or a prize)}

\label{lesson:JA-C457-oubo-suru}



\begin{quote}
Before the new one: say the Japanese for a job interview, then the Japanese for a curriculum vitae.
\end{quote}

\subsection*{You'll want to know: \ja{おうぼする}}
\textbf{\ja{おうぼする}} — \emph{ōbo suru} — \textquotedblleft{}to apply (for a job or a prize)\textquotedblright{}.

\begin{quote}
\textbf{\ja{かれは} \ja{みっつの} \ja{かいしゃに} \ja{おうぼしました。}} — \emph{kare wa mittsu no kaisha ni ōbo shimashita} — \textquotedblleft{}He applied to three companies.\textquotedblright{}
\end{quote}

Polite \textbf{\ja{おうぼします}}, \emph{ōboshimasu}; a \ja{する} verb. Past (polite): \textbf{\ja{がくせいが} \ja{おうぼしました。}}, \emph{gakusei ga ōbo shimashita}.

\begin{cousinweb}
\textbf{\ja{おうぼする}} (\emph{ōbo suru}) is Sino-Japanese (kango) \emph{ōbo} with the native verb \emph{suru}, \textquotedblleft{}to do\textquotedblright{}: two Chinese characters read with their on'yomi, readings borrowed from Chinese. They are \emph{ō}, \textquotedblleft{}answer\textquotedblright{} (Middle Chinese \emph{'ingH}) and \emph{bo}, \textquotedblleft{}recruit\textquotedblright{} (\emph{muH}) (KANJIDIC2; Middle Chinese in Baxter's transcription).
\end{cousinweb}

\begin{grammarlens}[title={What you've built}]
One more word for describing what you lived and felt.
\end{grammarlens}

\subsection*{Your turn}
\begin{itemize}
\raggedright
  \item \emph{Say it:} \emph{ōbo suru}
  \item \emph{Say it:} \emph{ōbo suru}, once more
  \item \emph{Say it:} \emph{kare wa mittsu no kaisha ni ōbo shimashita} — He applied to three companies.
  \item \emph{Recall:} say the Japanese for a job interview, then the Japanese for a curriculum vitae, then say \emph{ōbo suru} again
\end{itemize}

\begin{tcolorbox}[breakable,colback=teal!4,colframe=teal!35!black,title={Before you move on}]
What does \ja{おうぼする} mean? (\textquotedblleft{}To apply (for a job or a prize)\textquotedblright{}.) Say it once more.
\end{tcolorbox}
\section[\texorpdfstring{\ja{ぼしゅう} (boshū)}{ぼしゅう (boshū)}]{\ja{ぼしゅう} (boshū) — recruitment; a call for applicants}

\label{lesson:JA-C457-boshuu}



\begin{quote}
Before the new one: say the Japanese for a curriculum vitae, then the Japanese for to apply.
\end{quote}

\subsection*{You'll want to know: \ja{ぼしゅう}}
\textbf{\ja{ぼしゅう}} — \emph{boshū} — \textquotedblleft{}recruitment; a call for applicants\textquotedblright{}.

\begin{quote}
\textbf{\ja{みせで} \ja{てんいんを} \ja{ぼしゅうしています。}} — \emph{mise de ten'in o boshū shite imasu} — \textquotedblleft{}The shop is looking for staff.\textquotedblright{}
\end{quote}

\begin{cousinweb}
\textbf{\ja{ぼしゅう}} (\emph{boshū}) is Sino-Japanese (kango): two Chinese characters read with their on'yomi, readings borrowed from Chinese. They are \emph{bo}, \textquotedblleft{}recruit\textquotedblright{} (Middle Chinese \emph{muH}) and \emph{shū}, \textquotedblleft{}gather\textquotedblright{} (\emph{dzip}) (KANJIDIC2; Middle Chinese in Baxter's transcription).
\end{cousinweb}

\begin{grammarlens}[title={What you've built}]
One more word for describing what you lived and felt.
\end{grammarlens}

\subsection*{Your turn}
\begin{itemize}
\raggedright
  \item \emph{Say it:} \emph{boshū}
  \item \emph{Say it:} \emph{boshū}, once more
  \item \emph{Say it:} \emph{mise de ten'in o boshū shite imasu} — The shop is looking for staff.
  \item \emph{Recall:} say the Japanese for a curriculum vitae, then the Japanese for to apply, then say \emph{boshū} again
\end{itemize}

\begin{tcolorbox}[breakable,colback=teal!4,colframe=teal!35!black,title={Before you move on}]
What does \ja{ぼしゅう} mean? (\textquotedblleft{}Recruitment; a call for applicants\textquotedblright{}.) Say it once more.
\end{tcolorbox}
